% !TEX program = xelatex
% AI-effects 프로젝트 --- 한국 KSCO 8차 소분류 상위 30개(대화량, 대화량 × 절감 시간)와 연계 SOC 구성.
% 표 본문: results/table/26_summary.tex, 26_ksco3_top30_conv.tex, 26_ksco3_top30_saved.tex
%          (code/04_analysis_kor_ksco3_top30_soc.do가 생성)
% 전체 자료: results/table/26_kor_ksco3_top30_{conv,saved}_soc.xlsx (시트 ksco, soc)
% 관련: paper 19(한국 대화량), paper 21(연계표), paper 23(KSCO 결합), paper 25(활용도)
% 컴파일: xelatex (kotex 필요)
\documentclass[11pt]{article}

\usepackage{fontspec}
\usepackage{kotex}
\usepackage[a4paper,margin=2.2cm]{geometry}
\usepackage{longtable}
\usepackage{booktabs}
\usepackage{array}
\usepackage{amsmath}
\usepackage{xcolor}
\usepackage{hyperref}
\usepackage{enumitem}

\hypersetup{colorlinks=true,linkcolor=blue!50!black,urlcolor=blue!50!black}
\newcolumntype{L}[1]{>{\raggedright\arraybackslash}p{#1}}
\newcolumntype{R}[1]{>{\raggedleft\arraybackslash}p{#1}}
\setlength{\LTleft}{0pt}
\setlength{\LTright}{0pt}
\renewcommand{\arraystretch}{1.15}
\setlength{\tabcolsep}{4pt}

\title{한국 KSCO 소분류 상위 직업과 SOC 구성\\[2mm]
\large --- 업무용 AI 대화량과 절감 시간 상위 30개 소분류를 이루는 미국 직업 ---}
\author{AI-effects 프로젝트}
\date{\today}

\begin{document}
\maketitle

\begin{abstract}
\noindent 한국 업무용 AI 대화량(2026년 4$\cdot$5월 월평균, $K = 3$)을 KSCO 8차 소분류로 옮겼을 때 상위 30개 소분류를 두 기준으로 골랐다. 하나는 월 업무용 대화 수이고, 다른 하나는 대화 수와 대화당 절감 시간을 곱한 월 절감 시간이다. 그리고 각 소분류의 값이 어느 SOC 직업에서 왔는지 연계표 배분 몫으로 나누어 보았다. 주요 결과는 다음과 같다.
\begin{itemize}[leftmargin=*]
  \item 상위 30개 소분류가 소분류 전체 대화량의 81.2\%, 절감 시간의 84.1\%를 차지한다. 대화량 1위는 291 작가 및 언론 관련 전문가(월 4,126만 건), 절감 시간 1위는 222 컴퓨터 시스템 및 소프트웨어 전문가다.
  \item 소분류 하나의 값은 평균 10.7개 SOC에서 온다. 그러나 30개 중 14개는 SOC 하나가 절반 이상을 차지한다.
  \item SOC 대화량 1위인 15-1299(Computer Occupations, All Other)는 소분류 세 곳(225, 222, 223)으로 나뉘어 들어간다. 그래서 KSCO 순위에서는 단일 1위로 나타나지 않는다.
  \item paper 25의 활용도 1 1위인 213 생명 및 자연과학 관련 시험원은 대화량으로는 8위다. 그 대화의 34.5\%는 Data Scientists(15-2051)에서 온다.
\end{itemize}
\end{abstract}

\section{목적}

앞선 작업에서 한국 업무용 대화량을 SOC 단위(paper 19)와 KSCO 소분류 단위(paper 23, 25)로 만들었다. 두 순위는 서로 다르다. SOC 기준 1위는 15-1299지만, KSCO 기준 1위는 291 작가 및 언론 관련 전문가다. 이 노트는 KSCO 상위 소분류를 고른 뒤, 그 값을 이루는 SOC 직업을 펼쳐 이 차이가 어디서 오는지 보인다. 함께 만든 엑셀 파일에는 상위 30개 소분류와 연계된 모든 SOC가 들어 있다.

\section{자료와 계산}

\begin{itemize}[leftmargin=*]
  \item \textbf{SOC 값}: \texttt{kor\_soc6\_conv\_ai\_indicators}(paper 23)의 월별 업무용 대화 수 $c_s$($K = 3$)와 SOC 단계 절감 시간 $v_s = c_s \times \tau_s / 60$. 여기서 $\tau_s$는 AEI 전 세계 대화당 절감 시간(분)이다.
  \item \textbf{배분 몫}: SOC 2018--KSCO 8차 연계표(paper 21)의 $w^{\text{alloc}}(s,k)$를 소분류 $K$ 안에서 더한 값이다.
    \[
      w_3(s,K) = \sum_{k \in K} w^{\text{alloc}}(s,k), \qquad \sum_K w_3(s,K) = 1.
    \]
  \item \textbf{기여분}: SOC $s$가 소분류 $K$에 주는 대화 수는 $c_s \times w_3(s,K)$, 절감 시간은 $v_s \times w_3(s,K)$다. 월평균은 소분류 값과 같은 방식으로 냈다. 즉 그 소분류에 값이 있는 달 수로 나눈다. 따라서 기여분의 합은 paper 25의 소분류 값($C$, $S$)과 정확히 같다(do 파일에서 확인).
  \item \textbf{연계 SOC}: 소분류의 세분류 하나 이상과 연계표로 이어진 SOC 전부다. 한국 대화량이 공개되지 않은 SOC는 기여가 빈칸이다.
  \item \textbf{SOC 명칭}: BLS SOC 2018 공식 명칭이다(2026년 10월 18일 통일).
\end{itemize}

\section{결과}

\subsection{요약}

\begin{longtable}{L{8.5cm}R{3cm}R{3cm}}
\caption{상위 30개 소분류 요약}\label{tab:summary}\\
\toprule
 & 대화량 기준 & 절감 시간 기준 \\
\midrule
\endfirsthead
% 04_09_analysis_kor_ksco3_top30_soc.do가 생성. 직접 고치지 말 것.
값이 있는 소분류 수 & 129 & 129 \\
상위 30개 비중(\%) & 81.3 & 84.1 \\
상위 30개 평균 기여 SOC 수 & 10.7 & 10.9 \\
1위 SOC 기여 50\% 이상인 소분류 수 & 14 & 13 \\

\bottomrule
\end{longtable}

\noindent 주: 소분류 합계의 단위는 대화량이 만 건, 절감 시간이 만 시간(월평균)이다. ``평균 기여 SOC 수''는 상위 30개 소분류에서 한국 대화량이 공개된 연계 SOC 수의 평균이다.

\subsection{대화량 상위 30개 소분류}

\begin{longtable}{R{0.6cm}L{4.3cm}R{1.3cm}R{1.0cm}R{1.1cm}L{5.2cm}R{1.0cm}}
\caption{업무용 대화량 상위 30개 소분류와 기여 1위 SOC}\label{tab:conv}\\
\toprule
순위 & 소분류 & 월 대화 수(만 건) & 비중(\%) & 공개/연계 SOC & 기여 1위 SOC & 1위 비중(\%) \\
\midrule
\endfirsthead
\toprule
순위 & 소분류 & 월 대화 수(만 건) & 비중(\%) & 공개/연계 SOC & 기여 1위 SOC & 1위 비중(\%) \\
\midrule
\endhead
% 04_analysis_kor_ksco3_top30_soc.do가 생성. 직접 고치지 말 것.
1 & 291 작가 및 언론 관련 전문가 & 4,126 & 11.92 & 5/6 & 27-3041 Editors & 30.5 \\
2 & 222 컴퓨터 시스템 및 소프트웨어 전문가 & 3,444 & 9.95 & 11/11 & 15-1299 Computer Occupations, All Other & 24.2 \\
3 & 225 정보 시스템 및 웹 운영자 & 2,878 & 8.31 & 6/6 & 15-1299 Computer Occupations, All Other & 57.8 \\
4 & 372 비서 및 사무 지원 종사자 & 1,624 & 4.69 & 18/26 & 43-9022 Word Processors and Typists & 32.6 \\
5 & 283 상품 기획‧홍보 및 조사 전문가 & 1,453 & 4.20 & 9/9 & 27-3043 Writers and Authors & 57.3 \\
6 & 223 네트워크 및 정보 보안 전문가 & 1,328 & 3.84 & 7/8 & 15-1299 Computer Occupations, All Other & 62.6 \\
7 & 292 학예사‧사서 및 기록물 관리사 & 1,318 & 3.81 & 5/6 & 25-4022 Librarians and Media Collections Specialists & 76.3 \\
8 & 213 생명 및 자연과학 관련 시험원 & 964 & 2.79 & 14/24 & 15-2051 Data Scientists & 34.5 \\
9 & 332 통계‧데이터 관련 사무원 & 904 & 2.61 & 8/12 & 15-2051 Data Scientists & 36.8 \\
10 & 340 금융 사무 종사자 & 904 & 2.61 & 7/15 & 13-2071 Credit Counselors & 79.0 \\
11 & 293 연극‧영화 및 영상 전문가 & 840 & 2.43 & 9/20 & 27-3023 News Analysts, Reporters, and Journalists & 56.4 \\
12 & 211 생명 및 자연과학 관련 전문가 & 708 & 2.04 & 16/27 & 19-1029 Biological Scientists, All Other & 31.2 \\
13 & 282 금융 및 보험 전문가 & 687 & 1.98 & 15/19 & 13-2052 Personal Financial Advisors & 20.2 \\
14 & 264 문리‧기술 및 예능 강사 & 641 & 1.85 & 9/11 & 25-3041 Tutors & 39.8 \\
15 & 235 기계‧로봇공학 기술자 및 시험원 & 590 & 1.71 & 16/25 & 17-2141 Mechanical Engineers & 43.7 \\
16 & 269 기타 교육 전문가 & 489 & 1.41 & 39/43 & 25-9031 Instructional Coordinators & 50.8 \\
17 & 131 연구‧교육 및 법률 관련 관리자 & 479 & 1.39 & 8/9 & 11-9032 Education Administrators, Kindergarten through Secondary & 44.0 \\
18 & 371 총무 사무원 및 대학 행정조교 & 461 & 1.33 & 2/3 & 43-9061 Office Clerks, General & 75.5 \\
19 & 212 인문 및 사회과학 전문가 & 438 & 1.26 & 10/11 & 19-3011 Economists & 28.7 \\
20 & 295 디자이너 & 412 & 1.19 & 4/8 & 27-1024 Graphic Designers & 76.3 \\
21 & 224 데이터 전문가 & 380 & 1.10 & 2/2 & 15-1243 Database Architects & 87.6 \\
22 & 522 상품 대여 종사자 & 372 & 1.07 & 1/2 & 41-2021 Counter and Rental Clerks & 100.0 \\
23 & 239 기타 공학 전문가 및 관련 종사자 & 370 & 1.07 & 14/23 & 17-2199 Engineers, All Other & 35.1 \\
24 & 271 법률 전문가 & 359 & 1.04 & 7/11 & 23-1011 Lawyers & 71.8 \\
25 & 122 마케팅 및 광고‧홍보 관리자 & 350 & 1.01 & 5/5 & 11-2011 Advertising and Promotions Managers & 72.2 \\
26 & 261 대학교수 및 강사 & 349 & 1.01 & 35/37 & 25-9044 Teaching Assistants, Postsecondary & 24.9 \\
27 & 284 감정‧기술 영업 및 중개 관련 종사자 & 348 & 1.01 & 11/17 & 41-9031 Sales Engineers & 32.7 \\
28 & 510 영업 종사자 & 318 & 0.92 & 5/8 & 41-9011 Demonstrators and Product Promoters & 43.4 \\
29 & 121 행정 및 경영 지원 관리자 & 301 & 0.87 & 15/15 & 11-3131 Training and Development Managers & 30.6 \\
30 & 521 매장 판매 종사자 & 287 & 0.83 & 7/9 & 41-2011 Cashiers & 82.3 \\

\bottomrule
\end{longtable}

\begin{itemize}[leftmargin=*]
  \item \textbf{291 작가 및 언론 관련 전문가}가 1위다. 연계된 6개 SOC 중 5개가 공개되어 있고, 글쓰기 직업이 고르게 기여한다. Editors 30.5\%, Technical Writers 20.9\%, Writers and Authors 20.2\%, Interpreters and Translators 16.9\%, News Analysts, Reporters, and Journalists 11.5\%다. SOC 기준으로는 여러 직업에 흩어져 있던 대화가 한 소분류로 모인 경우다.
  \item \textbf{15-1299가 세 소분류의 1위 SOC다}. 225 정보 시스템 및 웹 운영자(57.8\%), 223 네트워크 및 정보 보안 전문가(62.6\%), 222 컴퓨터 시스템 및 소프트웨어 전문가(24.2\%)다. 연계표는 15-1299의 대화를 세분류 2251에 1/2, 2229에 1/4, 2232$\cdot$2239에 각 1/8씩 보낸다. 15-1299는 SOC 단계에서 월 약 3,300만 건으로 가장 크지만 이렇게 나뉘므로 KSCO에서는 단일 1위가 되지 않는다. 또 그 대화의 69\%는 Document Management Specialists(15-1299.03)의 과업이라, 정보 시스템 운영자 등으로 가는 배분은 실제 과업과 어긋날 수 있다(paper 23).
  \item \textbf{213 생명 및 자연과학 관련 시험원}은 8위(월 964만 건)다. 기여 1위가 Data Scientists(34.5\%)이고, 332 통계$\cdot$데이터 관련 사무원도 같다. 연계표가 Data Scientists를 ISCO 3314를 거쳐 두 세분류(2133, 3320)에 반씩 보내기 때문이다. paper 25에서 213이 활용도 1의 1위인 것은 이 대화가 취업자 1만 1천 명의 작은 소분류에 들어갔기 때문이다.
  \item \textbf{SOC 하나가 지배하는 소분류}도 많다. 30개 중 14개는 1위 SOC가 절반 이상이다. 522 상품 대여 종사자(Counter and Rental Clerks 100\%), 224 데이터 전문가(Database Architects 87.6\%), 521 매장 판매 종사자(Cashiers 82.3\%)가 그 예다. 이런 소분류의 값은 사실상 SOC 하나의 값이다.
\end{itemize}

\subsection{대화량 $\times$ 절감 시간 상위 30개 소분류}

\begin{longtable}{R{0.6cm}L{4.3cm}R{1.3cm}R{1.0cm}R{1.1cm}L{5.2cm}R{1.0cm}}
\caption{월 절감 시간 상위 30개 소분류와 기여 1위 SOC}\label{tab:saved}\\
\toprule
순위 & 소분류 & 월 절감 시간(만 시간) & 비중(\%) & 공개/연계 SOC & 기여 1위 SOC & 1위 비중(\%) \\
\midrule
\endfirsthead
\toprule
순위 & 소분류 & 월 절감 시간(만 시간) & 비중(\%) & 공개/연계 SOC & 기여 1위 SOC & 1위 비중(\%) \\
\midrule
\endhead
% 04_analysis_kor_ksco3_top30_soc.do가 생성. 직접 고치지 말 것.
1 & 222 컴퓨터 시스템 및 소프트웨어 전문가 & 18,263 & 12.21 & 11/11 & 15-1251 Computer Programmers & 26.2 \\
2 & 291 작가 및 언론 관련 전문가 & 15,668 & 10.47 & 5/6 & 27-3041 Editors & 32.3 \\
3 & 225 정보 시스템 및 웹 운영자 & 7,983 & 5.34 & 6/6 & 15-1299 Computer Occupations, All Other & 58.5 \\
4 & 372 비서 및 사무 지원 종사자 & 7,838 & 5.24 & 18/26 & 43-9022 Word Processors and Typists & 29.0 \\
5 & 283 상품 기획‧홍보 및 조사 전문가 & 7,700 & 5.15 & 9/9 & 27-3043 Writers and Authors & 52.3 \\
6 & 332 통계‧데이터 관련 사무원 & 6,027 & 4.03 & 8/12 & 15-2051 Data Scientists & 34.1 \\
7 & 223 네트워크 및 정보 보안 전문가 & 5,603 & 3.74 & 7/8 & 15-1299 Computer Occupations, All Other & 41.7 \\
8 & 213 생명 및 자연과학 관련 시험원 & 5,410 & 3.62 & 14/24 & 15-2051 Data Scientists & 38.0 \\
9 & 211 생명 및 자연과학 관련 전문가 & 4,321 & 2.89 & 16/27 & 19-1029 Biological Scientists, All Other & 39.1 \\
10 & 293 연극‧영화 및 영상 전문가 & 4,261 & 2.85 & 9/20 & 27-3023 News Analysts, Reporters, and Journalists & 46.9 \\
11 & 282 금융 및 보험 전문가 & 3,714 & 2.48 & 15/19 & 13-2051 Financial and Investment Analysts & 24.9 \\
12 & 292 학예사‧사서 및 기록물 관리사 & 3,421 & 2.29 & 5/6 & 25-4022 Librarians and Media Collections Specialists & 59.7 \\
13 & 264 문리‧기술 및 예능 강사 & 3,091 & 2.07 & 9/11 & 25-3041 Tutors & 54.7 \\
14 & 269 기타 교육 전문가 & 3,025 & 2.02 & 39/43 & 25-9031 Instructional Coordinators & 61.8 \\
15 & 224 데이터 전문가 & 2,723 & 1.82 & 2/2 & 15-1243 Database Architects & 91.5 \\
16 & 340 금융 사무 종사자 & 2,660 & 1.78 & 7/15 & 13-2071 Credit Counselors & 86.1 \\
17 & 235 기계‧로봇공학 기술자 및 시험원 & 2,542 & 1.70 & 16/25 & 17-2141 Mechanical Engineers & 39.1 \\
18 & 122 마케팅 및 광고‧홍보 관리자 & 2,246 & 1.50 & 5/5 & 11-2011 Advertising and Promotions Managers & 67.4 \\
19 & 131 연구‧교육 및 법률 관련 관리자 & 2,019 & 1.35 & 8/9 & 11-9032 Education Administrators, Kindergarten through Secondary & 40.4 \\
20 & 239 기타 공학 전문가 및 관련 종사자 & 1,990 & 1.33 & 14/23 & 17-2199 Engineers, All Other & 40.2 \\
21 & 121 행정 및 경영 지원 관리자 & 1,980 & 1.32 & 15/15 & 11-3131 Training and Development Managers & 36.2 \\
22 & 271 법률 전문가 & 1,873 & 1.25 & 7/11 & 23-1011 Lawyers & 66.9 \\
23 & 212 인문 및 사회과학 전문가 & 1,679 & 1.12 & 10/11 & 19-3011 Economists & 36.2 \\
24 & 261 대학교수 및 강사 & 1,669 & 1.12 & 35/37 & 25-1126 Philosophy and Religion Teachers, Postsecondary & 31.4 \\
25 & 295 디자이너 & 1,514 & 1.01 & 4/8 & 27-1024 Graphic Designers & 63.0 \\
26 & 371 총무 사무원 및 대학 행정조교 & 1,509 & 1.01 & 2/3 & 43-9061 Office Clerks, General & 66.3 \\
27 & 298 문화‧예술 관련 기획자 & 1,470 & 0.98 & 7/10 & 13-2099 Financial Specialists, All Other & 32.3 \\
28 & 141 건설‧전기 및 생산 관련 관리자 & 1,244 & 0.83 & 6/7 & 11-3051 Industrial Production Managers & 75.6 \\
29 & 234 전기‧전자공학 기술자 및 시험원 & 1,241 & 0.83 & 5/5 & 17-2061 Computer Hardware Engineers & 83.1 \\
30 & 236 소방‧방재‧산업안전 기술자 및 시험원 & 1,198 & 0.80 & 14/29 & 17-2199 Engineers, All Other & 44.6 \\

\bottomrule
\end{longtable}

\begin{itemize}[leftmargin=*]
  \item 상위권의 구성은 대화량 기준과 거의 같다. 30개 중 26개가 두 목록에 같이 있다. 대화량 목록에만 있는 것은 510 영업 종사자, 521 매장 판매 종사자, 522 상품 대여 종사자, 284 감정$\cdot$기술 영업 및 중개 관련 종사자다. 절감 시간 목록에만 있는 것은 298 문화$\cdot$예술 관련 기획자, 141 건설$\cdot$전기 및 생산 관련 관리자, 234 전기$\cdot$전자공학 기술자 및 시험원, 236 소방$\cdot$방재$\cdot$산업안전 기술자 및 시험원이다. 판매$\cdot$대여 직업은 대화당 절감 시간이 짧아 빠지고, 관리자$\cdot$공학 직업은 길어서 들어온다.
  \item \textbf{222 컴퓨터 시스템 및 소프트웨어 전문가}가 1위로 올라선다. 대화량은 2위지만 Computer Programmers(26.2\%) 등 대화당 절감 시간이 긴 SOC가 많아서다.
  \item 1위 SOC가 바뀌는 소분류도 있다. 222는 15-1299에서 Computer Programmers로, 282 금융 및 보험 전문가는 Personal Financial Advisors에서 Financial and Investment Analysts로, 261 대학교수 및 강사는 Teaching Assistants, Postsecondary에서 Philosophy and Religion Teachers, Postsecondary로 바뀐다.
\end{itemize}

\section{해석상 주의}

\begin{enumerate}[leftmargin=*]
  \item \textbf{배분 가정}: 소분류 값은 연계표의 배분 가중치에 크게 의존한다(paper 21). 15-1299나 Data Scientists처럼 큰 SOC가 어떻게 나뉘는지가 소분류 순위를 바꾼다. 연계 경로가 낯선 기여도 있다. 예를 들어 298 문화$\cdot$예술 관련 기획자의 절감 시간 1위가 13-2099 Financial Specialists, All Other인 것은 연계표 경로 때문이다.
  \item \textbf{AEI 직업의 의미}: AEI의 SOC는 사용자의 직업이 아니라 대화 과업이 속한 직업이다. 따라서 소분류 값은 ``그 소분류 종사자의 대화''가 아니라 ``그 소분류의 과업에 해당하는 대화''다.
  \item \textbf{절감 시간}: Claude가 대화별로 추정한 전 세계 평균이다(실측 아님). 수준보다 순위 비교에 적합하다(paper 25).
  \item \textbf{미공개 SOC}: 한국 대화량이 공개되지 않은 SOC는 기여가 없다. 연계 SOC 중 일부만 공개된 소분류(예: 372의 26개 중 18개)의 값은 하한에 가깝다.
  \item \textbf{$K$}: 대화 수는 $K$에 반비례할 뿐이라 순위와 비중은 $K$와 무관하다.
\end{enumerate}

\section{산출 파일}

\begin{itemize}[leftmargin=*]
  \item \texttt{code/04\_analysis\_kor\_ksco3\_top30\_soc.do}: 이 노트의 모든 표와 엑셀 파일을 만든다.
  \item \texttt{results/table/26\_kor\_ksco3\_top30\_conv\_soc.xlsx}: 대화량 상위 30개 소분류. 시트 \texttt{ksco}(소분류 요약), \texttt{soc}(소분류 $\times$ 연계 SOC 전부, 기여 대화 수$\cdot$비중$\cdot$절감 시간).
  \item \texttt{results/table/26\_kor\_ksco3\_top30\_saved\_soc.xlsx}: 절감 시간 상위 30개 소분류. 같은 구성.
  \item \texttt{results/table/26\_summary.tex}, \texttt{26\_ksco3\_top30\_conv.tex}, \texttt{26\_ksco3\_top30\_saved.tex}: 표~\ref{tab:summary}--\ref{tab:saved}.
\end{itemize}

\end{document}
